\section{Lecture 22 (November 26th)}
\begin{prop}
(Feynman rules in the momentum space) 
\begin{itemize}
\item[(i)] Construct all connected diagrams.
\item[(ii)] For each arrowed leg, we have the assign the Dirac propagator
\[\Big(\dfrac{-i}{\cancel{k}+m-i\epsilon }\Big)_{\alpha \beta }\]
\item[(iii)] For each wavy leg, in the Feynman gauge,
\[\dfrac{-i\eta _{\mu \nu }}{k^2-i\epsilon }\]
\item[(iv)] For each vertex, 
\[ie(\gamma ^{\mu })_{\beta \alpha }\cdot (2\pi )^{4}\delta ^{(4)}(k_1+k_2-k_3)\]
\item[(v)] Integrate over all the internal momenta.
\item[(vi)] Divide by the symmtry factor.
\item[(vii)] Consider the rearrangement of Grassmannian fields $\times (-1)^{n}$.
\item[(viii)] Sum over all diagrams.
\end{itemize}
\end{prop}
\vspace{2ex}
\begin{ex}
($e^{-}e^{+}\rightarrow e^{-}e^{+}$ Bhabha scattering) The correlation function of our interest is:
\[G^{(2;2;0)}_{c\;\alpha_1\alpha_2;\beta_1\beta_2}(x_1,x_2,y_1,y_2)=\langle T\psi _{\alpha_1}(x_1)\psi _{\alpha_2}(x_2)\bar{\psi }_{\beta_1}(y_1)\bar{\psi }_{\beta_2}(y_2)\]
The Fourier transform is given as
\[\tilde{G}^{(2;2;0)}_{c\;\alpha_1\alpha_2;\beta_1\beta_2}(p_1,p_2;q_1,q_2)\]
Let's first compute in the position space! There are two possible diagrams.
\[\begin{align*}
&G^{(2;2;0)}_{c\;\alpha_1\alpha_2;\beta_1\beta_2}(x_1,x_2;y_1,y_2)\\=&\int d^{4}z_1\,d^{4}z_2\,\Big[(-1)\Big(S_{F}(x_1-z_1)ie\gamma ^{\mu }S_{F}(z_1-y_1)\Big)D_{F\,\mu \nu }(z_1-z_2)\Big(S_{F}(x_2-z_2)ie\gamma ^{\nu }S_{F}(z_2-y_2)\Big)_{\alpha_2\beta_2}\\&+\Big(S_{F}(x_2-z_1)ie\gamma ^{\mu }S_{F}(z_1-y_1)\Big)_{\alpha_2\beta_1}D_{F\,\mu \nu }(z_1-z_2)\Big(S_{F}(x_1-z_2)ie\gamma ^{\nu }S_{F}(z_2-y_2)\Big)_{\alpha_1\beta_2}\Big]
\end{align*}\]
where $(m,n)=(1,0)$. In the momentum space,
\[\begin{align*}
&\tilde{G}^{(2;2;0)}_{c\;\alpha_1\alpha_2;\beta_1\beta_2}(-p',q;-p,q')\\
=&(-1)\int \dfrac{d^{4}k}{(2\pi )^{4}}\,\Big(\dfrac{1}{-\cancel{p}+m-i\epsilon }ie\gamma ^{\mu }(2\pi )^{4}\delta ^{(4)}(p+p'-k)\dfrac{1}{\cancel{p}+m-i\epsilon }\Big)_{\alpha_1,\beta_1}\\
&\times \dfrac{-i\eta _{\mu \nu }}{k^2-i\epsilon }\times \Big(\dfrac{1}{\cancel{q}+m-i\epsilon }ie\gamma ^{\nu }(2\pi )^{4}\delta ^{(4)}(k-q-q')\dfrac{1}{-\cancel{q'}+m-i\epsilon }\Big)_{\alpha_2 \beta_2 }\\
&+\Big[(-p',\alpha_1)\leftrightarrow (q_1,\alpha_2)\Big]\\
=&\delta ^{(4)}(p+p'-q-q')\Big[(-1)\Big(\cdot \Big)_{\alpha_1\alpha_2}\times \dfrac{-i\eta _{\mu \nu }}{(p+p')^2-i\epsilon }\Big(\cdot \Big)_{\alpha_2\beta_2}+\ldots \Big]
\end{align*}\]
\end{ex}
\vspace{2ex}
\begin{defi}
(Mandelstam variables for $2\rightarrow 2$ scattering) We define
\begin{itemize}
\item[(i)] $s=-(p+p')^2$, the square of the total intial momentum.
\item[(ii)] $t=-(p-q)^2$, the square of the difference of the initial and final momenta of the most similar particles.
\item[(iii)] $u=-(p-q')^2$, the interchangable variable with $t$ by swapping final momenta. 
\end{itemize}
Let's make some comments. Bhabha scattering consists of a $s$-channel and a $t$-channel. In the case of Moller scattering, there is a $t$-channel and $s$-channel.
\end{defi}
\vspace{2ex}
\begin{ex}
($e^{-}\gamma \rightarrow e^{-}\gamma $ Compton scattering)
\[\tilde{G}_{c\;\alpha;\beta;\mu \nu }^{(1;1;2)}(P;Q;K_1;K_2)=\tilde{G}^{(1;1;2)}_{c\;\alpha;\beta;\mu \nu }(q;-p;-K,K')w\]
The above is equal to
\[\begin{align*}
&(2\pi )^{4}\delta ^{(4)}(p+K-q-K')\times \\&\Big[\dfrac{-i\eta _{\mu \rho }}{K^2-i\epsilon }\dfrac{-i\eta _{\sigma \nu }}{K'^2-i\epsilon }\Big(\dfrac{-i}{\cancel{q}+m-i\epsilon }ie\gamma ^{\rho }\dfrac{-i}{(\cancel{p+K})+m-i\epsilon }ie\gamma ^{\rho }\dfrac{-i}{\cancel{p}+m-i\epsilon }\Big)_{\alpha \beta }\Big]\\
&+\Big((K,\mu )\leftrightarrow (-K',\nu )\Big)
\end{align*}\]

\end{ex}
\vspace{2ex}

